\begin{mistake}{2026-05-07}{03}

\begin{problemcontent}
设
\[
a_n=\int_0^1 \sin(x^n)\,\mathrm{d}x,\qquad b_n=\int_0^1 \sin^n x\,\mathrm{d}x\quad (n=1,2,\ldots),
\]
证明：
\begin{enumerate}
  \item[(I)] \(0\le b_n\le a_n\)；
  \item[(II)] \(\displaystyle \lim_{n\to\infty}a_n=\lim_{n\to\infty}b_n=0\)。
\end{enumerate}
\end{problemcontent}

\begin{solutioncontent}
先证 \(0\le b_n\le a_n\)。显然当 \(x\in[0,1]\) 时，\(\sin x\ge 0\)，故
\[
b_n=\int_0^1\sin^n x\,\mathrm{d}x\ge 0.
\]

下面证明 \(\sin^n x\le \sin(x^n)\)。设
\[
f(t)=\frac{\sin t}{t}\quad (t\in(0,1]).
\]
则
\[
f'(t)=\frac{t\cos t-\sin t}{t^2}
      =\frac{\cos t\,(t-\tan t)}{t^2}<0,
\]
因为 \(t\in(0,1]\subset(0,\frac{\pi}{2})\) 时 \(\cos t>0\)，且 \(\tan t>t\)。所以 \(f(t)=\frac{\sin t}{t}\) 在 \((0,1]\) 上单调递减。

当 \(x\in(0,1)\) 时，有 \(0<x^n\le x\le 1\)，由 \(f\) 单调递减得
\[
\frac{\sin(x^n)}{x^n}\ge \frac{\sin x}{x}.
\]
两边同乘 \(x^n>0\)，得
\[
\sin(x^n)\ge x^{n-1}\sin x.
\]
又由 \(0\le \sin x\le x\) 得
\[
\sin^{n-1}x\le x^{n-1},
\]
于是
\[
x^{n-1}\sin x\ge \sin^{n-1}x\cdot \sin x=\sin^n x.
\]
故
\[
\sin(x^n)\ge \sin^n x\qquad (0<x<1).
\]
在端点 \(x=0,1\) 处不影响积分比较，因此由定积分比较定理得
\[
\int_0^1\sin^n x\,\mathrm{d}x\le \int_0^1\sin(x^n)\,\mathrm{d}x,
\]
即 \(b_n\le a_n\)。综上，\(0\le b_n\le a_n\)。

再证极限。由 \(u\ge 0\) 时 \(0\le \sin u\le u\)，令 \(u=x^n\)，得
\[
0\le \sin(x^n)\le x^n\qquad (0\le x\le 1).
\]
积分得
\[
0\le a_n=\int_0^1\sin(x^n)\,\mathrm{d}x\le \int_0^1x^n\,\mathrm{d}x=\frac{1}{n+1}.
\]
由于 \(\displaystyle \lim_{n\to\infty}\frac{1}{n+1}=0\)，由夹逼定理得
\[
\lim_{n\to\infty}a_n=0.
\]
又由第一部分 \(0\le b_n\le a_n\)，再次由夹逼定理得
\[
\lim_{n\to\infty}b_n=0.
\]
\end{solutioncontent}

\mistakeknowledge{
\begin{itemize}
  \item \textbf{核心公式/定理}：定积分比较定理：若 \(f(x)\ge g(x)\) 于区间上成立，则 \(\int f\ge \int g\)。夹逼定理常用于证明含参积分极限为 \(0\)。
  \item \textbf{易错点}：不要用泰勒展开硬证全区间积分不等式；本题关键是先证被积函数不等式 \(\sin^n x\le \sin(x^n)\)。
  \item \textbf{关联知识}：函数 \(\frac{\sin t}{t}\) 在 \((0,1]\) 上单调递减，可由 \(\left(\frac{\sin t}{t}\right)'=\frac{\cos t(t-\tan t)}{t^2}<0\) 证明。
  \item \textbf{常用放缩}：当 \(u\ge 0\) 时，\(0\le \sin u\le u\)；因此 \(0\le a_n\le \int_0^1x^n\,\mathrm{d}x=\frac{1}{n+1}\).
\end{itemize}}

\end{mistake}
